\section{Quantum Foundations}

The base is local and deterministic. Quantum mechanics is nonlocal in the precise sense that it violates a bound any local model must respect. That is a sharp tension, and waving it away would misrepresent the model. This section states the tension plainly, then shows how the substrate's own holographic structure resolves it, and finally addresses two further questions, where the Born measure comes from and where the complex unit of quantum amplitudes comes from.

\subsection{The Bell tension}

A local, deterministic rule is a local hidden-variable theory, and every local hidden-variable theory obeys the Bell bound $\mathrm{CHSH} \le 2$. Brute-force enumeration over deterministic strategies confirms the maximum is $2$. Quantum mechanics reaches $2\sqrt{2} \approx 2.83$, the Tsirelson bound, verified at the optimal measurement angles. A local discrete base, taken literally as the whole story, therefore cannot reproduce quantum nonlocality. This is a genuine tension and not a free pass. Any discrete-substrate program has to confront it rather than assume it away.

\begin{result}[Bell bound versus Tsirelson]
A local deterministic rule obeys $\mathrm{CHSH} \le 2$, confirmed by enumeration over deterministic strategies. Quantum mechanics reaches $2\sqrt{2} \approx 2.83$ at the optimal angles \cite{pollard2026vibetest, bell1964, clauserhorne1969}. A literal local base cannot reach the quantum value.
\end{result}

\subsection{The resolution through holographic nonlocality}

The bulk rule is local, so the bulk obeys $\mathrm{CHSH} \le 2$. Physics, though, does not live in the bulk. It lives on the holographic boundary, and two boundary points are connected through the four-dimensional bulk by the bulk-tree propagator. The boundary theory is therefore nonlocal, with its correlations mediated through the bulk, and it is not a local hidden-variable model. A bulk-mediated boundary theory can violate Bell up to the quantum value while the rule that generates it stays strictly local. Nonlocality is emergent and holographic. The base remains local and discrete, and the quantum nonlocality is a property of the boundary, not of the substrate. This is consistent with the holographic structure the substrate already carries, and in fact it requires that structure, which is the same structure invoked in the gravity and bulk-cusp discussions.

\begin{proposition}[Holographic resolution of the Bell tension]
The local bulk obeys $\mathrm{CHSH} \le 2$, but the physical boundary theory is bulk-mediated and hence nonlocal, so it is not a local hidden-variable model and can violate Bell up to the quantum bound. Nonlocality is emergent and holographic while the base stays local \cite{maldacena1998, ryu2006, swingle2012}.
\end{proposition}

\subsection{The Born rule}

The $|\psi|^2$ measure appears as the natural measure of the oscillation amplitudes of the real discrete walk. That much is suggestive and consistent. The full account of measurement and collapse is the deeper open piece, and it is shared with every discrete and emergent-quantum approach rather than special to this one. It is listed here as open because it is open, not because it is solved.

\subsection{Where the complex unit comes from}

A natural worry is that the base is real and discrete, so the complex amplitude of quantum mechanics seems to have no place to come from. There is a clean answer through the Doi-Peliti map. The rule's three moves, create, annihilate, and hop, are all the same operator, the transfer of one unit of $S^z$ across an edge. The rule is therefore exactly a spin-$1$ $XX$ model, and its reversibility, which is detailed balance, is precisely what makes the Doi-Peliti Hamiltonian Hermitian and so gives a real, positive-energy quantum theory. In the spin-coherent-state path integral of that model the Berry-phase term is the source of the quantum $i$. The complex amplitude is not added. It is the Berry phase of the rule's own spin-exchange, and the imaginary unit is where the phase of the discrete walk lives.

\begin{proposition}[The $i$ as a spin Berry phase]
The rule's create, annihilate, and hop moves are one operator, the transfer of $S^z$ across an edge, so the rule is a spin-$1$ $XX$ model whose reversibility makes its Doi-Peliti Hamiltonian Hermitian. The Berry-phase term of the spin-coherent-state path integral is the source of the quantum $i$.
\end{proposition}


The section identifies a real requirement rather than papering over one. A discrete local base needs an emergent nonlocal boundary in order to be quantum, and the substrate supplies exactly that through its holography. There is a further pointer worth stating without overclaiming. Since the base is experiential under the monism, the quantum amplitudes are the texture of the underlying noticing, and the universe-as-mind is where the quantum layer ultimately sits. That is a candid pointer toward the rest of the program, not a derivation.
